\subsection{Mechanisms, assumptions, and implementation cost}
\owner{Member 2 integrates; each member completes their row.}
All methods use first-order gradients, but retain different optimizer states and perform
different scalar or coordinate-wise operations. Complexity should distinguish the gradient
calculation from the optimizer update. Record actual persistent state tensors rather than
only giving an $O(d)$ label.
\begin{table}[H]\centering\small
\caption{Method comparison and ownership. Complete cost and theorem details after implementation.}
\label{tab:methods}
\begin{tabularx}{\linewidth}{lXXl}\toprule
Pair & Main change & Assumptions and additional state & Writer\\\midrule
SGD / DoG & Dynamic scalar step & \tbd & (Member 1)\\
SGDM / SF-SGD & Averaging and interpolation & \tbd & (Member 2)\\
RMSProp / ADOPT-0 & Previous-moment normalization & \tbd & (Member 3)\\
Adam / AdEMAMix & Fast and slow gradient history & \tbd & (Member 4)\\\bottomrule
\end{tabularx}\end{table}
Each implementation will be checked on a small deterministic problem or fixed gradient
sequence against an independent calculation or trusted reference. Record discrepancies,
version choices, numerical tolerances and any intentional deviations.
